% Non-Monotonic Optimum in Prefix-Token Perturbation of Language Model Reasoning
% NeurIPS 2026 Workshop Submission — DRAFT
%
% Sections marked [LOCKED] are stable and can be finalized.
% Sections marked [PENDING] await experimental data.
% Sections marked [PLACEHOLDER] need substantial revision.

\documentclass{article}
\usepackage[utf8]{inputenc}
\usepackage{amsmath,amssymb}
\usepackage{booktabs}
\usepackage{graphicx}
\usepackage{hyperref}
\usepackage{xcolor}
\usepackage{subcaption}

% NeurIPS workshop format (if needed, swap for neurips_2026.sty)
\usepackage[margin=1in]{geometry}

\newcommand{\pending}[1]{\textcolor{red}{[PENDING: #1]}}
\newcommand{\placeholder}[1]{\textcolor{orange}{[PLACEHOLDER: #1]}}

\title{Random Prefix Tokens Reveal a Sharp Non-Monotonic Optimum\\in Language Model Arithmetic Reasoning}
% Title is [PENDING] per Codex — finalize after Shi comparison

\author{Devansh}
\date{March 2026}

\begin{document}
\maketitle

% ============================================================================
\begin{abstract}
Prepending random embedding-scale tokens to the input of a 4-bit quantized
language model (Qwen3-4B) produces a sharp, non-monotonic improvement in
chain-of-thought arithmetic reasoning. Two random tokens raise mean
accuracy from 32\% to 51.6\% (+19.6 percentage points) across ten
independent directions, while one token yields 42.7\% and three or eight
tokens revert to $\sim$44\%. The effect decomposes into think-mode gating
(+8pp) and perturbation-specific optimization (+12pp). Different
perturbation directions solve \emph{different task subsets}: ten two-token
directions collectively cover 100\% of tasks (25/25), compared to 80\%
for three-token and 92\% for eight-token perturbations. The three-token
oracle set is a strict subset of the two-token oracle, indicating that
over-perturbation contracts the reachable task space. These results
establish embedding-space perturbation as a training-free intervention
that activates latent reasoning capabilities in small language models.
\end{abstract}

% ============================================================================
\section{Introduction} % [LOCKED — stable claims only]

Large language models possess reasoning capabilities that are not
uniformly accessible through standard prompting. Recent work has shown
that semantically meaningless prefix tokens---repeated punctuation
\cite{shi2025}, pause tokens \cite{goyal2024}, and trailing whitespace
\cite{salinas2024}---can measurably alter model behavior, sometimes
improving reasoning performance.

We study a related but distinct regime: \emph{continuous} perturbation
in embedding space. Rather than prepending discrete tokens from the
vocabulary, we inject random vectors at the model's native embedding
scale as soft prompt prefixes. This eliminates semantic content entirely,
isolating the effect of perturbation energy and position from token
identity.

Our key findings on Qwen3-4B (4-bit quantized) with 25 nested arithmetic
tasks are:

\begin{enumerate}
    \item \textbf{Sharp non-monotonic optimum.} Two random prefix tokens
    raise mean accuracy from 32\% to 51.6\% across ten directions;
    one token yields 42.7\%, and three or eight tokens revert to
    $\sim$44\%.

    \item \textbf{Complete oracle coverage.} Ten two-token directions
    collectively solve 100\% of tasks (25/25), compared to 80\% (20/25)
    for three-token and 92\% (23/25) for eight-token perturbations.
    The three-token oracle set is a strict subset of the two-token
    oracle: five tasks solvable at two tokens become frozen at three.
    At matched budget ($k = 3$), two tokens achieve 77\% oracle coverage
    versus 60\% for three tokens and 68\% for eight tokens.

    \item \textbf{Heterogeneous task selectivity.}
    Different perturbation directions solve different task subsets,
    but per-direction solve counts vary normally (std = 1.87,
    $p = 0.66$ vs.\ heterogeneous iid null). The diversity is in
    \emph{which} tasks are solved, not in \emph{how many}.

    \item \textbf{Think-mode gating.} Any perturbation activates
    Qwen3's reasoning mode (16\% $\to$ 100\%), but random noise
    contributes $\sim$1.5$\times$ more accuracy gain than mode
    activation alone.
\end{enumerate}

All experiments use greedy decoding (temperature = 0), making the
observed variability across perturbation directions a form of
\emph{deterministic chaos} rather than sampling noise.

% ============================================================================
\section{Related Work} % [LOCKED — core refs stable]

\paragraph{Discrete prefix perturbations.}
Shi et al.~\cite{shi2025} demonstrate that repeated punctuation tokens
(e.g., ``/'' and ``?'') improve reasoning accuracy by 1--5\% across
multiple models, attributing the effect to MLP activation redistribution.
Our work studies the continuous embedding-space regime, finding
substantially larger effects (+19.6pp mean at the optimum), direction
independence, and task-selective redistribution---phenomena not observed
in discrete token space.
\pending{Add Shi head-to-head comparison when discrete token experiment completes.}

\paragraph{Pause tokens.}
Goyal et al.~\cite{goyal2024} show that training with learned pause
tokens improves downstream performance. Our effect requires no training:
random embeddings at inference time suffice.

\paragraph{Extra-token expressivity.}
London \& Nagarajan~\cite{london2025} prove that additional tokens
increase transformer expressivity. Our results provide an empirical
instantiation: two tokens at the correct energy scale unlock reasoning
that the model cannot access with standard prompting.

\paragraph{Deterministic chaos in transformers.}
Li et al.~\cite{li2025} introduce quasi-Lyapunov exponents for
language models, establishing that small perturbations to hidden states
grow exponentially through transformer layers. Our invariance-length
analysis (Section~\ref{sec:chaos}) provides a complementary behavioral
measure of this phenomenon under greedy decoding.

\paragraph{Stochastic resonance.}
The non-monotonic dose-response is reminiscent of stochastic resonance
\cite{gammaitoni1998}, where noise improves signal detection in nonlinear
systems. However, our statistical tests do not confirm the formal
inverted-U relationship at the task level (Section~\ref{sec:limitations}),
so we use the more conservative framing of ``non-monotonic optimum.''

% ============================================================================
\section{Methods} % [LOCKED]

\subsection{Model and Inference}

We use Qwen3-4B~\cite{qwen2025} with 4-bit GPTQ quantization,
running on a single NVIDIA RTX 5090 Laptop GPU (24GB VRAM). All
generation uses greedy decoding (temperature = 0, \texttt{do\_sample =
False}) with \texttt{max\_new\_tokens = 1024} and thinking mode enabled.

\subsection{Task Suite}

We use 25 nested arithmetic expressions at ``sweet-spot'' difficulty
(3-step expressions with intermediate results in the hundreds to
thousands range, 32\% baseline accuracy). Tasks are deterministically
generated and fixed across all conditions.

\subsection{Perturbation Protocol}

We inject $k$ random embedding vectors as soft prompt tokens before
the user input. Each vector is sampled from $\mathcal{N}(0, I)$ and
scaled to match the model's embedding table RMS norm
($\sigma_{\text{emb}} \approx 0.022$), calibrated via
\texttt{auto\_calibrate()} over 1000 vocabulary samples.

\subsection{Conditions}

\begin{table}[h]
\centering
\small
\begin{tabular}{llp{5cm}}
\toprule
\textbf{Condition} & \textbf{Tokens} & \textbf{Description} \\
\midrule
Baseline & 0 & No prefix tokens \\
Zero embedding & 8 & All-zero vectors \\
Mean embedding & 8 & Mean of embedding table \\
Force-think & 0 & Forced \texttt{<think>} token, no noise \\
Random noise & 1, 2, 3, 8 & Random $\mathcal{N}(0,I)$, RMS-matched \\
\pending{Discrete tokens} & \pending{1, 2, 8} & \pending{Repeated `/', `?' tokens (Shi comparison)} \\
\bottomrule
\end{tabular}
\caption{Experimental conditions. All use greedy decoding.}
\label{tab:conditions}
\end{table}

\subsection{Scoring and Data Integrity}

Correctness is determined by \texttt{verify\_answer()}: the last integer
in the generated response must exactly match the expected answer. All
experiments with the current codebase store the extracted answer alongside
the full response (up to 2000 characters) for independent verification.

% ============================================================================
\section{Results} % [LOCKED sections noted; PENDING sections noted]

\subsection{Non-Monotonic Dose-Response} % [LOCKED]
\label{sec:dose}

\begin{table}[h]
\centering
\begin{tabular}{cccccc}
\toprule
\textbf{Tokens} & \textbf{Accuracy} & \textbf{Std} & \textbf{$n$} & \textbf{$\Delta$ vs baseline} \\
\midrule
0 & 32.0\% & --- & 1 & --- \\
1 & 42.7\% & 1.9\% & 3 & +10.7pp \\
\textbf{2} & \textbf{51.6\%} & \textbf{7.9\%} & \textbf{10} & \textbf{+19.6pp} \\
3 & 44.0\% & 5.3\% & 10 & +12.0pp \\
8 & 44.4\% & 7.0\% & 10 & +12.4pp \\
\bottomrule
\end{tabular}
\caption{Dose-response: accuracy vs.\ number of random prefix tokens
($n$ = number of independent directions tested).
Two tokens is the sharp optimum. Three and eight tokens revert to
$\sim$44\%.}
\label{tab:dose}
\end{table}

\begin{figure}[h]
\centering
\includegraphics[width=0.7\linewidth]{experiments/fig_dose_response.pdf}
\caption{Non-monotonic dose-response curve. Error bars show standard
error of the mean across $n$ perturbation directions.}
\label{fig:dose}
\end{figure}

The relationship between token count and accuracy is sharply non-monotonic.
Two random tokens produce the highest mean accuracy (51.6\% $\pm$ 7.9\%
across ten independent directions). Adding a third token drops accuracy to
$\sim$44\%, matching the eight-token level. The three-token oracle set
(20/25) is a strict subset of the two-token oracle (25/25), indicating
that over-perturbation contracts the reachable task space rather than
merely redistributing it.

\subsection{Force-Think Decomposition} % [LOCKED]
\label{sec:decomposition}

\begin{table}[h]
\centering
\begin{tabular}{lccc}
\toprule
\textbf{Condition} & \textbf{Think Rate} & \textbf{Accuracy} & \textbf{Component} \\
\midrule
Baseline & 16\% & 32.0\% & --- \\
Force-think (no noise) & 100\% & 40.0\% & +8pp (mode gate) \\
2-tok random noise & 100\% & 51.6\% & +11.6pp (perturbation) \\
\bottomrule
\end{tabular}
\caption{Two-component decomposition. Think-mode activation alone adds
+8pp; random noise contributes an additional +11.6pp beyond think mode
(mean across 10 directions).}
\label{tab:decomposition}
\end{table}

The perturbation effect has two separable components. First, any
perturbation activates Qwen3's reasoning mode (think rate: 16\% $\to$
100\%). Second, the noise itself produces an additional +11.6pp gain
beyond what think-mode activation provides (mean across 10 directions).
Noise contributes $\sim$1.5$\times$ more than mode gating at the
two-token optimum.

\subsection{Oracle Efficiency} % [LOCKED]
\label{sec:oracle}

Because different perturbation directions solve different task subsets,
running multiple perturbations and taking the oracle (any-correct)
provides coverage well above individual accuracy. Figure~\ref{fig:oracle}
shows coverage as a function of budget $k$ for each condition.

\begin{figure}[h]
\centering
\includegraphics[width=0.8\linewidth]{experiments/fig_oracle_coverage.pdf}
\caption{Oracle coverage vs.\ perturbation budget. Two-token (green)
reaches 100\% coverage at $k = 10$ (zero frozen tasks), dominates at
every matched budget, and its oracle set is a strict superset of the
three-token set (red, 80\% at $k = 10$). Eight-token (blue) reaches
92\% at $k = 10$.}
\label{fig:oracle}
\end{figure}

The non-monotonic optimum extends beyond per-direction accuracy to
coverage efficiency: five tasks solvable at two tokens become frozen
at three tokens, contracting the reachable task set rather than merely
redistributing it. Two-token perturbation alone covers all 25 tasks;
combined oracle across all conditions is therefore 25/25 = 100\%.

Oracle coverage grows more slowly than the independence null (positive
task correlation, $z = -1.35$ for 2-tok vs.\ $z = -8.25$ for 8-tok),
indicating that latents fail on correlated hard tasks. The two-token
condition has the least redundancy, closest to independence.

\subsection{Task Selectivity Without Equalization}
\label{sec:equalization}

An initial observation at $n = 3$ two-token directions suggested
solve-count equalization (all directions solving exactly 15/25 tasks,
$p = 0.031$ vs.\ heterogeneous iid null). This did not replicate:
at $n = 10$ directions, solve counts range from 9 to 15
(std = 1.87, $p = 0.66$), fully consistent with independent sampling
at heterogeneous per-task solve rates. The initial pattern was
small-sample noise. The robust finding is \emph{task selectivity}:
directions solve different subsets, enabling oracle gains, but the
total number solved per direction varies normally.

We test equalization by Monte Carlo simulation: given the observed
per-task solve rates, we simulate 100{,}000 draws under a
heterogeneous iid Bernoulli null (each task independently solved at
its observed rate) and compare the observed standard deviation of
per-direction solve counts against this null distribution.

We decompose tasks at each perturbation level into three categories:
\emph{unanimous} (solved by all sampled directions), \emph{frozen}
(solved by none), and \emph{sensitive} (solved by some but not all).
These categories are probabilistic, not absolute---additional directions
may solve ``frozen'' tasks or fail ``unanimous'' ones.

\begin{table}[h]
\centering
\textbf{Panel A: Matched budget ($k = 3$ directions)}

\medskip
\begin{tabular}{lcccc}
\toprule
\textbf{Condition} & \textbf{Unan.} & \textbf{Frozen} & \textbf{Sens.} & \textbf{Mean acc.} \\
\midrule
1-tok & 6.0 & 8.0 & 11.0 & 10.7/25 \\
\textbf{2-tok}$^\dagger$ & \textbf{7.6} & \textbf{5.8} & \textbf{11.6} & \textbf{12.9/25} \\
3-tok$^\dagger$ & 7.3 & 9.9 & 7.8 & 11.0/25 \\
8-tok$^\dagger$ & 6.0 & 7.9 & 11.2 & 11.1/25 \\
\bottomrule
\end{tabular}

\medskip
\textbf{Panel B: Full observed directions (with iid $p$-values)}

\medskip
\begin{tabular}{lccccccc}
\toprule
\textbf{Condition} & $k$ & \textbf{Unan.} & \textbf{Frozen} & \textbf{Sens.} & \textbf{Mean/25} & \textbf{Std} & $p_{\text{iid}}$ \\
\midrule
1-tok & 3 & 6 & 8 & 11 & 10.7 & 0.47 & 0.23 \\
\textbf{2-tok} & \textbf{10} & \textbf{6} & \textbf{0} & \textbf{19} & \textbf{12.9} & \textbf{1.87} & \textbf{0.659} \\
3-tok & 10 & 5 & 5 & 15 & 11.0 & 1.33 & 0.335 \\
8-tok & 10 & 3 & 2 & 20 & 11.1 & 1.76 & 0.586 \\
\bottomrule
\end{tabular}
\caption{Task decomposition by perturbation magnitude.
\textbf{Panel A}: At matched budget ($k = 3$), two tokens minimize the
frozen set and maximize per-direction accuracy.
$^\dagger$Values are averaged over all $\binom{n}{3}$ subsets
(120 combinations each for $n = 10$ conditions).
\textbf{Panel B}: Full observed directions with Monte Carlo
$p$-values against heterogeneous iid null. No condition shows
significant variance suppression; the diversity is in task identity,
not solve count.}
\label{tab:equalization}
\end{table}

\subsection{Deterministic Chaos} % [LOCKED]
\label{sec:chaos}

Under greedy decoding, different random prefix vectors produce
byte-identical text for an ``invariance length'' that decreases with
perturbation energy:

\begin{table}[h]
\centering
\begin{tabular}{ccc}
\toprule
\textbf{Tokens} & \textbf{Mean Invariance} & \textbf{Oracle} \\
\midrule
1 & 349 chars & 68\% \\
2 & 323 chars & 100\% \\
8 & 93 chars & 92\% \\
\bottomrule
\end{tabular}
\caption{Invariance length (byte-identical prefix before divergence)
under greedy decoding. This is deterministic chaos, not sampling noise.}
\label{tab:invariance}
\end{table}

The two-token sweet spot balances long enough invariance to preserve
reasoning quality with short enough invariance for task-subset diversity.

\subsection{Controls} % [LOCKED]

\begin{itemize}
    \item \textbf{Zero embedding} (+4pp): Sequence extension alone
    provides minimal benefit.
    \item \textbf{Mean embedding} (+4pp): Identical to zero---no
    advantage from the ``average token.''
    \item \textbf{No-think mode}: Perturbation has no effect when
    chain-of-thought is disabled. CoT is the mediating mechanism.
    \item \textbf{Timing confound resolved}: The highest-accuracy
    2-tok directions achieve 60\% (15/25) at baseline-identical mean
    timing (73.6s), ruling out extra computation as the mechanism.
\end{itemize}

\subsection{Task-Specific Resonance Windows} % [LOCKED]

Individual tasks exhibit condition-dependent resonance windows:
\begin{itemize}
    \item \texttt{nest\_005} (answer = 8360): solvable at 1-tok and
    2-tok but frozen at 3-tok and 8-tok
    \item \texttt{nest\_008} (answer = 7278): solvable \emph{only}
    at 2-tok (the sole task uniquely rescued by a single condition)
    \item \texttt{nest\_021} (answer = 5639): solvable at 2-tok and
    8-tok but frozen at 1-tok and 3-tok
\end{itemize}
More noise can \emph{prevent} solving certain tasks while enabling
others. At $n = 10$ directions, the 2-tok condition has zero frozen
tasks, while 3-tok freezes 5 and 8-tok freezes 2.

\subsection{Shi et al.\ Comparison} % [PENDING]
\label{sec:shi}

\pending{Head-to-head comparison of continuous (random embedding) vs.\
discrete (repeated punctuation) prefix tokens at matched token counts.
Experiment queued.}

% ============================================================================
\section{Discussion} % [PLACEHOLDER — refine after mechanism probes]

\subsection{Two-Stage Gating Model}

The structural decomposition suggests a two-stage model:
\begin{enumerate}
    \item \textbf{Magnitude sets the task-level threshold.}
    Perturbation energy determines which tasks are \emph{unanimous}
    (solved by all directions), \emph{frozen} (solved by none), or
    \emph{sensitive} (solved by some). At $n = 10$, two tokens freeze
    zero tasks, three tokens freeze 5, and eight tokens freeze 2.
    The three-token frozen set is a strict superset of the two-token
    frozen set: over-perturbation contracts reachability.

    \item \textbf{Direction selects within the sensitive set.}
    Each random direction solves a different subset of the sensitive
    tasks. Per-direction solve counts vary normally across all
    conditions (all $p > 0.3$ vs.\ heterogeneous iid null), so the
    diversity is in task identity, not in total solve count.
\end{enumerate}

The two-token optimum emerges because it maximizes the reachable task
set (25/25 at $n = 10$, vs.\ 20/25 for three tokens) while maintaining
the highest mean per-direction accuracy (12.9/25 = 51.6\%).

\subsection{Perturbation as Policy Switch}

The prefix perturbation does not make the model ``smarter''; it
switches the model's output policy from a formal presentation mode
(structured LaTeX, enumerated steps, truncation before computation)
to an exploratory computation mode (informal, stream-of-consciousness,
actual arithmetic). This policy switch helps tasks where the default
mode fails but causes overthinking on tasks the model already handles
efficiently.

\subsection{Why Two Tokens?}

At $n = 10$ directions, two tokens have zero frozen tasks, compared to
5 at three tokens and 2 at eight tokens. At matched budget ($k = 3$),
two tokens freeze 5.8 tasks on average versus 9.9 for three tokens
and 7.9 for eight tokens. Two tokens thus maximize the reachable task
set while maintaining the highest per-direction accuracy
(mean 12.9/25 = 51.6\%). We hypothesize that two tokens provide the
minimum perturbation energy needed to cross the mode-gating threshold
while staying below the energy level that freezes marginal tasks.
At one token, the perturbation is marginal; at three or more, the
model's reasoning trajectory diverges too early, exhausting the token
budget before reaching an answer.

\pending{Think-gate probe results will test this hypothesis directly.}

\subsection{Implications for Practice}

The training-free nature of this intervention suggests a practical
``perturbation ensemble'' strategy: run a small model with $k$ random
two-token prefixes and take the majority vote or oracle. Ten two-token
directions achieve 100\% coverage (25/25 tasks), while ten three-token
directions reach only 80\% and add no tasks beyond the two-token
oracle. Beyond the two-token optimum, additional perturbation is
deleterious: it lowers per-direction accuracy and contracts the
reachable task set, freezing tasks that are solvable at two tokens.

% ============================================================================
\section{Limitations} % [LOCKED]
\label{sec:limitations}

\begin{itemize}
    \item \textbf{Single model.} All experiments use Qwen3-4B (4-bit).
    The effect may not generalize to other architectures or scales.

    \item \textbf{Single task domain.} Only nested arithmetic
    expressions were tested. Generalization to reasoning, coding, or
    language tasks is unknown.
    \pending{Word problem replication may address this.}

    \item \textbf{Small $n$.} 25 tasks with 10 perturbation directions at
    the key 2-token condition. Per-task mechanistic analyses remain
    statistically limited.

    \item \textbf{4-bit quantization.} Quantization amplifies chaos
    (reduced numerical precision increases sensitivity to perturbation).
    Full-precision experiments may show different dose-response curves.

    \item \textbf{Stochastic resonance framing.} The delta-accuracy
    Spearman correlation with answer magnitude is not significant
    ($r = -0.295$, $p = 0.15$), and the logistic regression quadratic
    term is not significant ($p = 0.148$). The stochastic resonance
    analogy is suggestive but not statistically confirmed.

    \item \textbf{Token budget interaction.} Regressions on easy tasks
    (large-answer category: 100\% $\to$ 78\%) suggest the effect
    interacts with the 1024-token generation budget.
\end{itemize}

% ============================================================================
\section{Conclusion}

Random prefix tokens at the embedding scale reveal a sharp non-monotonic
optimum in small-model arithmetic reasoning. Two tokens produce a
mean +19.6 percentage point improvement over baseline, decomposable into
think-mode gating (+8pp) and perturbation-specific optimization (+12pp).
Different random directions solve different task subsets, enabling
complete oracle coverage (100\% from 10 two-token directions, compared
to 80\% for three-token and 92\% for eight-token). The three-token
oracle set is a strict subset of the two-token oracle, with five tasks
becoming frozen at three tokens. Per-direction solve counts vary
normally ($p = 0.66$ vs.\ iid null); the diversity is in task identity,
not total count.
These findings suggest that small language models possess latent
reasoning capabilities that can be unlocked through minimal,
training-free embedding-space perturbation.

% ============================================================================
\bibliographystyle{plain}
% References (inline for now; move to .bib for submission)
\begin{thebibliography}{99}

\bibitem{shi2025}
Shi, Z., et al.
\newblock Meaningless tokens, meaningful gains: Repeated tokens improve LLM reasoning.
\newblock \emph{arXiv:2510.01032}, 2025.

\bibitem{goyal2024}
Goyal, S., et al.
\newblock Think before you speak: Training language models with pause tokens.
\newblock \emph{ICLR}, 2024.

\bibitem{salinas2024}
Salinas, D. and Morstatter, F.
\newblock The butterfly effect of altering prompts: How small changes and jailbreaks affect large language model performance.
\newblock \emph{arXiv:2401.03729}, 2024.

\bibitem{london2025}
London, B. and Nagarajan, V.
\newblock Provably more expressive transformers via extra tokens.
\newblock \emph{NeurIPS}, 2025.

\bibitem{li2025}
Li, Z., et al.
\newblock Quasi-Lyapunov exponent for language models.
\newblock \emph{arXiv}, 2025.

\bibitem{gammaitoni1998}
Gammaitoni, L., et al.
\newblock Stochastic resonance.
\newblock \emph{Reviews of Modern Physics}, 70(1):223, 1998.

\bibitem{qwen2025}
Qwen Team.
\newblock Qwen3 technical report.
\newblock 2025.

\bibitem{xu2025}
Xu, Z., et al.
\newblock Causal amplification in the residual stream.
\newblock \emph{arXiv}, 2025.

\end{thebibliography}

\end{document}
